\newpage\sloppy
\section{Cross-Cutting Security \& Governance}
ระบบรักษาความปลอดภัยครอบคลุมทั้งระบบ ในแนวตั้ง เพื่อให้มั่นใจในความน่าเชื่อถือของเครือข่ายโลจิสติกส์ สถาปัตยกรรมจึงกำหนดให้ต้องปฏิบัติตามกรอบการรักษาความปลอดภัยระดับสากลอย่างเคร่งครัด\par

\begin{itemize}
	\item \textbf{Centralized Role \& Matrix}: แก้ปัญหาเดิมที่สิทธิ์การใช้งาน (Roles) กระจัดกระจายตามโมดูลต่างๆ โดยนำมารวมศูนย์ไว้ที่เดียว เพื่อให้ตรวจสอบและแก้ไขได้ง่าย\par
	\item \textbf{Audit Logging}: มีการใช้โมดูล \textbf{sso-idm-metabase} สำหรับวิเคราะห์บันทึกการตรวจสอบ (Audit Logs) และรายงานการปฏิบัติตามข้อกำหนด\par
	\item \textbf{SOC Integration}: มีแผนในการเชื่อมต่อ Log ของระบบไปยังศูนย์ปฏิบัติการความปลอดภัย (SOC) เพื่อการเฝ้าระวังแบบ Real-time\par
	\item \textbf{PDPA}: ระบบถูกออกแบบมาเพื่อตอบโจทย์ของ ICT Audit ในเรื่องการปฏิบัติตามพระราชบัญญัติคุ้มครองข้อมูลส่วนบุคคล (PDPA) ของประเทศไทย\par
	\item \textbf{Centralized Revocation}: ช่วยให้Chief Security Officer (CSO) สามารถเพิกถอนสิทธิ์การเข้าถึงของผู้รับเหมาภายนอกหรือพนักงานได้จากจุดเดียวทันที\par
\end{itemize}

\begin{itemize}
	\item ขั้นตอนการพัฒนาและดูแลรักษาได้รับการรับรองมาตรฐาน ISO/IEC 27001: เพื่อให้มั่นใจว่าระบบการจัดการความปลอดภัยของข้อมูล (ISMS) ที่ครอบคลุมนั้นถูกนำไปใช้ในกระบวนการ SSO และการจัดการข้อมูลผู้ใช้
	\item ได้รับการรับรองมาตรฐาน ISO/IEC 27017: เพื่อให้มั่นใจว่ามีการควบคุมความปลอดภัยที่เข้มงวดเฉพาะ สำหรับการจัดหาและการใช้บริการคลาวด์ (โครงสร้างพื้นฐาน AWS)
\end{itemize}
\fussy
